\section{Results}

Table~\ref{tab:design-space-results} front-loads the main quantitative summary
for judging, so this section only records the interpretation. Across the sweep,
wider input buses and more aggressive neuron parallelism reduce latency and
raise throughput, but they also push harder on timing closure and memory
distribution. That trend is consistent with the structure of the RTL: the
compute side can consume more work per step only if the replay path, banked
RAMs, and neuron processors all scale with it.

The implementation strategy for that tradeoff is visible in the architecture
figures. Configuration writes are staged and fanned out through explicit
pipeline registers, while the compute path uses replay buffering and separate
read-control logic to keep activation availability decoupled from local RAM
access sequencing. Final synthesis and implementation artifacts are stored
under \texttt{openflex/}, especially the CSV summaries in
\texttt{openflex/csv\_runs/}. The submitted operating point should therefore be
the row that best balances throughput gain against LUT, register, BRAM, and
timing cost once the final resource numbers are locked.
